\documentclass[../algebraIII_notes.tex]{subfiles}
\begin{document}
\section{Aula 13 - 26 de Maio, 2026}
\subsection{Motivações}
\begin{itemize}
	\item Teorema Fundamental da Teoria de Galois.
\end{itemize}
\subsection{Teorema Fundamental da Teoria de Galois}
Se \(\mathbb{K}/\mathbb{F}\) for uma extensão de Galois, denotamos o grupo de Galois da extensão por
\[
	\mathrm{Gal}(\mathbb{K}/\mathbb{F}) \coloneqq \mathrm{Aut}_{\mathbb{F}}(\mathbb{K},
\]
tal que, pelos teoremas anteriores,
\[
	\mathrm{ord}(\mathrm{Gal}(\mathbb{K}/\mathbb{F})) = [\mathbb{K}:\mathbb{F}].
\]
Hoje, enunciamos a provamos o
\hypertarget{fundamental_galois}{
	\begin{theorem*}[Teorema Fundamental da Teoria de Galois]
		Para corpos \(\mathbb{K}\) e \(\mathbb{F}\) com \(\mathbb{K}/\mathbb{F}\), temos:
		\begin{itemize}
			\item[1)] O morfismo
			      \[
				      \{H < \mathrm{Gal}(\mathbb{K}/\mathbb{F})\}\stackrel{\mathrm{Inv}(\cdot )}\longrightarrow  \{\mathbb{F}\subseteq \mathbb{M}\subseteq \mathbb{M}:\; \mathfrak{M} \text{ subcorpo}\}
			      \]
			      é uma bijeção com inversa
			      \[
				      \{\mathbb{F}\subseteq \mathbb{M}\subseteq \mathbb{K}:\; \mathbb{M} \text{ corpo}\}\stackrel{\mathrm{Gal}(\mathbb{K}/(\cdot ))}\longrightarrow  \{H\subseteq \mathrm{Gal}(\mathbb{K}/\mathbb{F})\};
			      \]
			\item[2)] Vale que \(H_1\) é um subgrupo de \(H_2\) se, e somente se, \(\mathrm{Inv}(H_2)\subseteq \mathrm{Inv}(H_1)\);
			\item[3)] Há uma relação entre índice de subgrupos e grau de extensão de \(\mathbb{H}_1/\mathbb{H}_2\):
			      \[
				      [H_2: H_1]= [\mathrm{Inv}(H_1):\mathrm{Inv}(H_2)];
			      \]
			\item[4)] H é um subgrupo normal de \(\mathrm{Gal}(\mathbb{K}/\mathbb{F})\) se, e somente sem a extensão \(\mathrm{Inv}(H)/\mathbb{F}\) for uma extensão normal. Neste caso, há um isomorfismo de grupos
			      \[
				      \mathrm{Gal}(\mathrm{Inv}(H)/\mathbb{F})\cong \frac{\mathrm{Gal}(\mathbb{K}/\mathbb{F})}{H};
			      \]
			\item[5)] Para toda \(\phi \in \mathrm{Gal}(\mathbb{K}/\mathbb{F})\) e todo \(H < \mathrm{Gal}(\mathbb{K}/\mathbb{F})\),
			      \[
				      \mathrm{Inv}(\phi H \phi^{-1}) = \phi (\mathrm{Inv}(H)).
			      \]
			      Em particular, se \(\mathbb{F}\subseteq \mathbb{M}\subseteq \mathbb{M}\), então
			      \[
				      \mathrm{Gal}(\mathbb{K}/\phi (\mathbb{M})) = \phi  \mathrm{Gal}(\mathbb{K}/\mathbb{M})\phi^{-1}.
			      \]
		\end{itemize}
	\end{theorem*}}

\begin{proof*}
	Vamos provar alguns itens, mas não todos, de modo explícito.

	(1):  para começar, observe duas coisas: se \(\mathbb{K}/\mathbb{F}\) for de Galois, então
	\[
		\mathrm{ord}(\mathrm{Gal}(\mathbb{K}/\mathbb{F}))< +\infty,
	\]
	e, para todo corpo \(\mathbb{M}\) intermediário da torre \(\mathbb{F}\subseteq \mathbb{M}\subseteq \mathbb{K}\), \(\mathbb{K}/\mathbb{M}\) é Galois, afinal basta observar que se \(\mathbb{K}/\mathbb{F}\) é finita, então \(\mathbb{K}/\mathbb{M}\) é finita, pois \(\mathbb{K}/\mathbb{F}\) é Galois se, e só se, \(\mathbb{K}\) é o corpo de decomposição de um polinômio \(f(x)\subseteq \mathbb{F}[x]\) separável, e por isso \(\mathbb{K}\) é corpo de decomposição também de \(f(x)\in \mathbb{M}[x]\).
	Daí, o item (1) segue ao considerar \(H < \mathrm{Gal}(\mathbb{K}/\mathbb{F})\), de forma que \(\mathrm{ord}(H)\) é finita e, consequentemente, por um teorema que foi consequência do \hyperlink{artin_lemma}{\textit{Lema de Artin}}, segue o lado
	\[
		\mathrm{Gal}(\mathbb{K}/\mathrm{Inv}(H)) = H.
	\]
	Por outro lado, se \(\mathbb{F}\subseteq \mathbb{M}\subseteq \mathbb{K}\), então já mostramos que
	\[
		\mathrm{Inv}(\mathrm{Gal}(\mathbb{K}/\mathbb{M})) = \mathbb{M}.
	\]

	(2): se \(H_1 < H_2\), é sempre verdade que \(\mathrm{Inv}(H_2)\subseteq \mathrm{Inv}(H_1)\), e que, então,
	\[
		\mathrm{Gal}(\mathbb{K}/\mathrm{Inv}(H_1)) \subseteq \mathrm{Gal}(\mathbb{K}/\mathrm{Inv}(H_2)).
	\]
	No entanto, sendo \(\mathbb{K}/\mathbb{F}\) Galois, deve ocorrer
	\[
		\mathrm{Gal}(\mathbb{K}/\mathrm{Inv}(H_1)) = H_1 \quad\&\quad \mathrm{Gal}(\mathbb{K}/\mathrm{Inv}(H_2))=H_2.
	\]
	Logo,
	\[
		H_1 < H_2 \Longleftrightarrow \mathrm{Inv}(H_2) < \mathrm{Inv}(H_1).
	\]

	(3): sabemos que, se H for subgrupo de \(\mathrm{Gal}(\mathbb{K}/\mathbb{F})\), então
	\[
		\mathrm{ord}(\mathrm{Gal}(\mathbb{K}/\mathrm{Inv}(H))) = [\mathbb{K}:\mathrm{Inv}(H)].
	\]
	Sejam \(H_1 < H_2\) dois grupos finitos; com isto,
	\begin{align*}
		\mathrm{ord}(H_2) = [H_2:\overbrace{1}^{\mathclap{\substack{\text{Índice do} \\\text{grupo trivial} }}}] &= [H_2:H_1][H_1:1]\\
		 & = [H_2:H_1]\mathrm{ord}(H_1).
	\end{align*}
	Além disso, para qualquer índice de H< vale
	\[
		\mathrm{ord}(H_{i}) = [\mathbb{K}:\mathrm{Inv}(H_1)]= \mathrm{ord}(\mathrm{Gal}(\mathbb{K}/\mathrm{Inv}(H_i))),
	\]
	e pela multiplicidade dos graus,
	\[
		[\mathbb{K}:\mathrm{Inv}(H_2)] = [\mathbb{K}:\mathrm{Inv}(H_1)][\mathrm{Inv}(H_1):\mathrm{Inv}(H_2)],
	\]
	implicando em
	\[
		\mathrm{ord}(H_2)=\mathrm{ord}(H_1)[\mathrm{Inv}(H_1):\mathrm{Inv}(H_2)].
	\]
	Juntando ambas, obtemos:
	\[
		[\mathrm{Inv}(H_1):\mathrm{Inv}(H_2)]\frac{\mathrm{ord}(H_2)}{\mathrm{ord}(H_1)} = [H_2:H_1].
	\]

	(4): Seja H um subgrupo normal de \(\mathrm{Gal}(\mathbb{K}/\mathbb{F}),\) ou seja, suponha que, para todo \(\phi \in \mathrm{Gal}(\mathbb{K}/\mathbb{F}),\) existe um elemento de H que é conjugado do original: \(\phi H \phi^{-1}=H\).
	Pelo exercício logo abaixo, podemos definir um homomorfismo de grupo
	\begin{align*}
		\tau : & \mathrm{Gal}(\mathbb{K}/\mathbb{F})\rightarrow \mathrm{Aut}(\mathrm{Inv}(H)/\mathbb{F}) \\
		       & \phi\longmapsto \phi|_{\mathrm{Inv}(H)}
	\end{align*}
	Usando o teorema da aula anterior e os exercícios abaixo, podemos concluir que \(\mathrm{Inv}(H)/\mathbb{F}\) é de Galois pois
	\[
		\mathbb{K}/\mathbb{F} \text{ é Galois } \Longleftrightarrow \exists G \subseteq \mathrm{Gal}(\mathbb{K}/\mathbb{F})
	\]
	é finita e \(\mathbb{F} = \mathrm{Inv}(G)\), e basta trocar G pela imagem de \(\tau \). Por sobrejetividade de \(\tau \) e pelo \hyperlink{isomorphism_theorem}{\textit{Primeiro Teorema do Isomorfismo}} para grupos, segue que
	\[
		\mathrm{Gal}(\mathrm{Inv}(H)/\mathbb{F})\cong \frac{\mathrm{Gal}(\mathbb{K/\mathbb{F}})}{H}.
	\]

	Por outro lado, Seja \(\mathbb{M}\) um corpo intermediário de \(\mathbb{F}\subseteq \mathbb{M}\subseteq \mathbb{K}\), tal que \(\mathbb{M}/\mathbb{F}\) é normal. Resta provar que \(H\coloneqq \mathrm{Gal}(\mathbb{K}/\mathbb{H})\) é normal, ou seja, que para toda \(\phi \in \mathrm{Gal}(\mathbb{K}/\mathbb{F})\), temos \(\phi H\phi^{-1}=H\).

	Com efeito, seja \(\mathbb{M} = \mathbb{F}[\alpha_1, \dotsc , \alpha_{m}]\), onde cada \(\alpha_{i}\) é algébrico sobre \(\mathbb{F}\); se \(\phi \in \mathrm{Gal}(\mathbb{K}/\mathbb{F})\) e aplicarmos \(\phi \) ao elemento \(\alpha_{i}\), obtemos uma raiz do polinômio minimal de \(\alpha_{i}\), pois os coeficientes dele são elementos de \(\mathbb{F}\), que são fixados. Logo, dada uma raiz, aplicar \(\varphi \) dá outra, para todo \(\alpha_{i}\), e por \(\mathbb{M}/\mathbb{F}\) ser normal, também vale sempre \(\varphi (\alpha_{i})\in \mathbb{M}\). Portanto,
	\[
		\varphi M = M (**).
	\]

\end{proof*}
\begin{exr}
	Mostre que H é um subgrupo normal de \(\mathrm{Gal}(\mathbb{K}/\mathbb{F})\) se, e somente se, \(\mathrm{Inv}(H)\) é \(\mathrm{Gal}(\mathbb{K}/\mathbb{F})\)-invariante, ou seja, se para todo \(\phi \in \mathrm{Gal}(\mathbb{K}/\mathbb{F}),\) vale
	\[
		\phi \mathrm{Inv}(H) = \mathrm{Inv}(H).
	\]
	Consequentemente, que podemos definir um homomorfismo de grupo
	\begin{align*}
		\tau : & \mathrm{Gal}(\mathbb{K}/\mathbb{F})\rightarrow \mathrm{Aut}(\mathrm{Inv}(H)/\mathbb{F}) \\
		       & \phi\longmapsto \phi|_{\mathrm{Inv}(H)}.
	\end{align*}
	Mostre que \(\tau \) é sobrejetora.

	Além disso, mostre que \(\mathrm{ker}(\tau )=H\), onde
	\[
		\mathrm{ker}(\tau ) = \{\phi \in \mathrm{Gal}(\mathbb{K}/\mathbb{F}):\; \tau (\phi ) = \mathrm{Id}_{\mathrm{Inv(H)}}\}
	\]
	e que o diagrama abaixo comuta
	\begin{center}
		\begin{tikzpicture}[
				observed/.style = {rectangle, thick, text centered, draw, text width = 6em},
				latent/.style = {ellipse, thick, draw, text centered, text width = 6em},
				error/.style ={circle, thick, draw, text centered},
				confounding/.style = {rectangle, thick, text centered, draw, text width = 6em, minimum width = 5.5in},
				outcome/.style = {rectangle, thick, draw, text centered, minimum height = 3.5in, text width = 6em},
			]

			\node(T) at (-2,1){\(\mathrm{Gal}(\mathbb{K}/\mathbb{F})\)};
			\node(BL) at (-2,-1){\(\displaystyle\frac{\mathrm{Gal}(\mathbb{K}/\mathbb{F})}{H}\)};
			\node(BR) at (2,1){\(\mathrm{Aut}(\mathrm{Inv}(H)/\mathbb{F})\)};

			\draw[Arrow](T)--node[midway, left] {\(\pi \)}(BL);
			\draw[Arrow](T)--node[midway, above] {\(\tau \)}(BR);
			\draw[Arrow, dashed](BL)--node[midway, below] {}(BR);

		\end{tikzpicture}
	\end{center}
	Equivalentemente, mostre que
	\[
		\frac{\mathrm{Gal}(\mathbb{K}/\mathbb{F})}{H}\cong \mathrm{Im}(\tau ).
	\]
\end{exr}
\begin{exr}
	Seja
	\[
		\mathrm{Inv}_{\mathrm{Inv}(H)}\coloneqq \{a\in \mathrm{Inv}(H): \phi (a) = a,\; \forall \phi \in \mathrm{Im}(\tau )\}.
	\]
	Mostre que
	\[
		\mathrm{Inv}_{\mathrm{Inv}(H)}(\mathrm{Im}(\tau )) = \mathbb{F}.
	\]
\end{exr}
\begin{exr}
	Mostre que (**) implica que \(\mathrm{Gal}(\mathbb{K}/\mathbb{F})\eqqcolon\) é normal.
\end{exr}
\begin{exr}
	Mostre o item (5).
\end{exr}

Adiante, veremos um exemplo simples e um exemplo mais complexo, que irá requerer introduzir e trabalhar por um tempo com as chamadas classes ciclotrômicas.

\end{document}
